\textbf{Phần 1. Trắc nghiệm lựa chọn}
\setcounterpageref{ex}{0}

\begin{ex} %[4D11V60]
Muốn cộng hai phân số cùng mẫu số, ta
\choice
{Cộng tử số với tử số, cộng mẫu số với mẫu số}
{\True Cộng hai tử số với nhau và giữ nguyên mẫu số}
{Nhân hai tử số với nhau}
{Quy đồng tử số}
\loigiai{
Cộng hai tử số với nhau và giữ nguyên mẫu số.\\
\textbf{Đáp án của câu hỏi là B}
}
\end{ex}

\begin{ex} %[4D11V60]
Kết quả của phép tính $\dfrac{2}{7} + \dfrac{3}{7}$ là
\choice
{$\dfrac{5}{14}$}
{$\dfrac{6}{7}$}
{$\dfrac{1}{7}$}
{\True $\dfrac{5}{7}$}
\loigiai{
$\dfrac{2 + 3}{7} = \dfrac{5}{7}$.\\
\textbf{Đáp án của câu hỏi là D}
}
\end{ex}

\begin{ex} %[4D11V60]
Kết quả của phép tính $\dfrac{1}{3} + \dfrac{1}{6}$ là
\choice
{\True $\dfrac{1}{2}$}
{$\dfrac{2}{9}$}
{$\dfrac{3}{6}$}
{$\dfrac{2}{6}$}
\loigiai{
$\dfrac{2}{6} + \dfrac{1}{6} = \dfrac{3}{6} = \dfrac{1}{2}$.\\
\textbf{Đáp án của câu hỏi là A}
}
\end{ex}

\begin{ex} %[4D11V60]
Kết quả của phép tính $\dfrac{3}{4} + \dfrac{2}{5}$ là
\choice
{$\dfrac{5}{9}$}
{$\dfrac{11}{20}$}
{\True $\dfrac{23}{20}$}
{$\dfrac{6}{20}$}
\loigiai{
$\dfrac{15}{20} + \dfrac{8}{20} = \dfrac{23}{20}$.\\
\textbf{Đáp án của câu hỏi là C}
}
\end{ex}

\begin{ex} %[4D11V60]
Kết quả của phép tính $1 + \dfrac{2}{5}$ là
\choice
{$\dfrac{3}{5}$}
{\True $\dfrac{7}{5}$}
{$\dfrac{2}{5}$}
{$\dfrac{8}{5}$}
\loigiai{
$\dfrac{5}{5} + \dfrac{2}{5} = \dfrac{7}{5}$.\\
\textbf{Đáp án của câu hỏi là B}
}
\end{ex}

\begin{ex} %[4D11V60]
Tính thuận tiện: $\dfrac{1}{4} + \dfrac{2}{7} + \dfrac{3}{4}$.
\choice
{$1$}
{$\dfrac{8}{7}$}
{$2$}
{\True $\dfrac{9}{7}$}
\loigiai{
$\left(\dfrac{1}{4} + \dfrac{3}{4}\right) + \dfrac{2}{7} = 1 + \dfrac{2}{7} = \dfrac{9}{7}$.\\
\textbf{Đáp án của câu hỏi là D}
}
\end{ex}

\begin{ex} %[4D11V60]
Một người thợ ngày thứ nhất dệt được $\dfrac{2}{7}$ tấm vải, ngày thứ hai dệt được $\dfrac{3}{7}$ tấm vải. Cả hai ngày người đó dệt được số phần tấm vải là
\choice
{\True $\dfrac{5}{7}$ tấm vải}
{$\dfrac{5}{14}$ tấm vải}
{$\dfrac{6}{7}$ tấm vải}
{$1$ tấm vải}
\loigiai{
$\dfrac{2}{7} + \dfrac{3}{7} = \dfrac{5}{7}$ tấm vải.\\
\textbf{Đáp án của câu hỏi là A}
}
\end{ex}

\begin{ex} %[4D11V60]
Một vòi nước giờ thứ nhất chảy được $\dfrac{1}{3}$ bể nước, giờ thứ hai chảy được $\dfrac{2}{5}$ bể nước. Cả hai giờ vòi chảy được số phần bể nước là
\choice
{$\dfrac{3}{8}$ bể}
{$\dfrac{7}{15}$ bể}
{\True $\dfrac{11}{15}$ bể}
{$\dfrac{2}{15}$ bể}
\loigiai{
$\dfrac{1}{3} + \dfrac{2}{5} = \dfrac{5 + 6}{15} = \dfrac{11}{15}$ bể.\\
\textbf{Đáp án của câu hỏi là C}
}
\end{ex}

\begin{ex} %[4D11V60]
Tìm $x$ biết: $x - \dfrac{1}{4} = \dfrac{5}{8}$.
\choice
{$\dfrac{6}{8}$}
{\True $\dfrac{7}{8}$}
{$\dfrac{3}{8}$}
{$1$}
\loigiai{
$x = \dfrac{5}{8} + \dfrac{2}{8} = \dfrac{7}{8}$.\\
\textbf{Đáp án của câu hỏi là B}
}
\end{ex}

\begin{ex} %[4D11V60]
Tính giá trị biểu thức: $\dfrac{1}{2} + \dfrac{1}{3} + \dfrac{1}{6}$.
\choice
{$\dfrac{5}{6}$}
{$\dfrac{7}{6}$}
{$\dfrac{3}{11}$}
{\True $1$}
\loigiai{
$\dfrac{3}{6} + \dfrac{2}{6} + \dfrac{1}{6} = \dfrac{6}{6} = 1$.\\
\textbf{Đáp án của câu hỏi là D}
}
\end{ex}

\begin{ex} %[4D11V60]
Một hình tam giác có độ dài 3 cạnh lần lượt là $\dfrac{3}{5}\text{ m}, \dfrac{4}{5}\text{ m}, \dfrac{2}{5}\text{ m}$. Chu vi tam giác đó là
\choice
{\True $\dfrac{9}{5}\text{ m}$}
{$\dfrac{9}{15}\text{ m}$}
{$2\text{ m}$}
{$\dfrac{7}{5}\text{ m}$}
\loigiai{
$\dfrac{3 + 4 + 2}{5} = \dfrac{9}{5}\text{ m}$.\\
\textbf{Đáp án của câu hỏi là A}
}
\end{ex}

\begin{ex} %[4D11V60]
Tính tổng: $S = \dfrac{1}{1 \times 2} + \dfrac{1}{2 \times 3} + \dfrac{1}{3 \times 4} + \dfrac{1}{4 \times 5}$.
\choice
{$\dfrac{3}{5}$}
{$\dfrac{5}{6}$}
{\True $\dfrac{4}{5}$}
{$1$}
\loigiai{
$S = 1 - \dfrac{1}{2} + \dfrac{1}{2} - \dfrac{1}{3} + \dots - \dfrac{1}{5} = 1 - \dfrac{1}{5} = \dfrac{4}{5}$.\\
\textbf{Đáp án của câu hỏi là C}
}
\end{ex}

\textbf{Phần 2. Trắc nghiệm đúng sai}
\setcounterpageref{ex}{0}

\begin{ex} %[4D11V60]
Xét tính đúng sai của các phép tính cộng phân số:
\choiceTF
{\True $\dfrac{1}{5} + \dfrac{2}{5} = \dfrac{3}{5}$.}
{$\dfrac{1}{2} + \dfrac{1}{3} = \dfrac{2}{5}$.}
{\True $\dfrac{1}{2} + \dfrac{1}{3} = \dfrac{5}{6}$.}
{\True $\dfrac{3}{8} + \dfrac{5}{8} = 1$.}
\loigiai{
\begin{itemize}
    \item Cùng mẫu số: $1/5 + 2/5 = 3/5$. Đúng.
    \item $1/2 + 1/3 = 5/6 \neq 2/5$. Ý 2 Sai, Ý 3 Đúng.
    \item $3/8 + 5/8 = 8/8 = 1$. Đúng.
\end{itemize}
\begin{itemize}
    \item a) \textbf{Đúng}.
    \item b) \textbf{Sai}.
    \item c) \textbf{Đúng}.
    \item d) \textbf{Đúng}.
\end{itemize}
}
\end{ex}

\begin{ex} %[4D11V60]
Hai vòi nước cùng chảy vào một bể cạn: Vòi 1 chảy được $\dfrac{2}{7}$ bể, vòi 2 chảy được $\dfrac{3}{7}$ bể. Xét tính đúng sai:
\choiceTF
{\True Cả hai vòi cùng chảy được $\dfrac{5}{7}$ bể nước.}
{\True Vòi 2 chảy được nhiều hơn vòi 1 là $\dfrac{1}{7}$ bể nước.}
{Sau khi hai vòi chảy thì bể đã đầy nước.}
{\True Còn lại $\dfrac{2}{7}$ bể chưa có nước.}
\loigiai{
\begin{itemize}
    \item Tổng: $2/7 + 3/7 = 5/7$ bể. Đúng.
    \item Hiệu: $3/7 - 2/7 = 1/7$ bể. Đúng.
    \item $5/7 < 1$ nên chưa đầy. Sai.
    \item Còn lại: $1 - 5/7 = 2/7$ bể. Đúng.
\end{itemize}
\begin{itemize}
    \item a) \textbf{Đúng}.
    \item b) \textbf{Đúng}.
    \item c) \textbf{Sai}.
    \item d) \textbf{Đúng}.
\end{itemize}
}
\end{ex}

\textbf{Phần 3. Trả lời ngắn}
\setcounterpageref{ex}{0}

\begin{ex} %[4D11V60]
Tính tổng: $\dfrac{4}{9} + \dfrac{2}{9}$ (rút gọn tối giản dạng a/b).
\shortans{2/3}
\loigiai{
\textbf{Đáp số:} $2/3$\\
$6/9 = 2/3$.
}
\end{ex}

\begin{ex} %[4D11V60]
Tính tổng: $\dfrac{1}{4} + \dfrac{3}{8}$ (dạng a/b).
\shortans{5/8}
\loigiai{
\textbf{Đáp số:} $5/8$\\
$2/8 + 3/8 = 5/8$.
}
\end{ex}

\begin{ex} %[4D11V60]
Tính giá trị: $1 + \dfrac{3}{4}$ (dạng a/b).
\shortans{7/4}
\loigiai{
\textbf{Đáp số:} $7/4$\\
$7/4$.
}
\end{ex}

\begin{ex} %[4D11V60]
Tìm $x$ biết: $x - \dfrac{2}{5} = \dfrac{1}{2}$ (dạng a/b).
\shortans{9/10}
\loigiai{
\textbf{Đáp số:} $9/10$\\
$x = 1/2 + 2/5 = 5/10 + 4/10 = 9/10$.
}
\end{ex}

\begin{ex} %[4D11V60]
Tính thuận tiện: $\dfrac{2}{9} + \dfrac{5}{7} + \dfrac{7}{9}$ (dạng a/b).
\shortans{12/7}
\loigiai{
\textbf{Đáp số:} $12/7$\\
$1 + 5/7 = 12/7$.
}
\end{ex}

\begin{ex} %[4D11V60]
Tính tổng: $\dfrac{1}{2} + \dfrac{1}{4} + \dfrac{1}{8} + \dfrac{1}{16}$ (dạng a/b).
\shortans{15/16}
\loigiai{
\textbf{Đáp số:} $15/16$\\
$8/16 + 4/16 + 2/16 + 1/16 = 15/16$.
}
\end{ex}
